\documentclass[11pt,a4paper,dvipdfmx]{article} 		% for platex and uplatex
\input{pieces/form00_header} % pieces
\input{pieces/kakenhi7} % pieces
\input{pieces/form01_header} % pieces
\input{pieces/form02_2022_header} % pieces
\input{pieces/hook3} % pieces
%#Name: ilr-3
\input{pieces/form04_jsps_eng_headers} % pieces
\input{pieces/form04_ilr-3_header} % pieces
% ===== Global definitions for the Kakenhi form ======================
%　基本情報
%
%------ 研究課題名  -------------------------------------------
\newcommand{\研究課題名}{象の卵}

%----- 研究機関名と研究代表者の氏名-----------------------
\newcommand{\研究機関名}{逢坂大学}
\newcommand{\研究代表者氏名}{湯川秀樹}
\newcommand{\me}{\underline{\underline{H.~Yukawa}}} 
\input{pieces/inst_general} % pieces
\input{pieces/inst_ilr-3} % pieces
% user07_header
% ===== my favorite packages ====================================
% ここに、自分の使いたいパッケージを宣言して下さい。
\usepackage{wrapfig}
%\usepackage{amssymb}
%\usepackage{mb}
%\DeclareGraphicsRule{.tif}{png}{.png}{`convert #1 `dirname #1`/`basename #1 .tif`.png}
\usepackage{lineno}

% ===== my personal definitions ==================================
% ここに、自分のよく使う記号などを定義して下さい。
\newcommand{\klpionn}{K_L \to \pi^0 \nu \overline{\nu}}
\newcommand{\kppipnn}{K^+ \to \pi^+ \nu \overline{\nu}}

% ----- 業績リスト用 -------------
\newcommand{\paper}[6]{%
	% paper{title}{authors}{journal}{vol}{pages}{year}
	\item ``#1'', #2, #3 {\bf #4}, #5 (#6).			% お好みに合わせて変えてください。
}

\newcommand{\etal}{\textit{et al.\ }}
\newcommand{\ca}[1]{*#1}	% corresponding author;   \ca{\yukawa}  みたいにして使う
\newcommand{\invitedtalk}{招待講演}

\newcommand{\yukawa}{H.~Yukawa}					% no underline
%\newcommand{\yukawa}{\underline{\underline{H.~Yukawa}}}	% with 2 underlines
\newcommand{\tomonaga}{S.~Tomonaga}

\newcommand{\prl}{Phys.\ Rev.\ Lett.\ }		% よく使う雑誌も定義すると楽

% ===== 欄外メモ ==================
\newcommand{\memo}[1]{\marginpar{#1}}
%\renewcommand{\memo}[1]{}	% 全てのメモを表示させないようにするには、行頭の"%"を消す

%%\input{../../pieces/user07_header}	% expand
%\input{../../sample/simple/contents}	% skip
\input{pieces/hook5} % pieces

\begin{document}
\input{pieces/hook7} % pieces
%#Split: 01_PI_requirements  
%#PieceName: p01_PI_requirements
\input{pieces/p01_PI_requirements_00}
\section{Application Requirements that PI must meet 
}
%　　　　＜＜最大　2ページ＞＞

% s13_pi_requirements_ilr_en
%begin piRequirements ＝＝＝＝＝＝＝＝＝＝＝＝＝＝＝＝＝＝＝＝
\JSPSInstructions	% <-- 留意事項。これは消すか、コメントアウトしてください。

	The requirements on PI are as follows.
	\begin{itemize}
		\item The PI should have an indomitable spirit to
			search for elephant eggs even under strong 
			doubts and criticisms by orthodox biologists. 
		\item To lead an international group, the PI should be
			fluent in multiple languages.
		\item To communicate with elephants, the PI should be
			fluent in elephant's languages, both 
			in African and Indian dialects.
		\item Most importantly, the PI should have a pure heart of a child.
	\end{itemize}
	
	\input{blahblah}
%end piRequirements ＝＝＝＝＝＝＝＝＝＝＝＝＝＝＝＝＝＝＝＝

\input{pieces/p01_PI_requirements_01}

%#Split: members_info % keep

% Principal Investigator　＝＝＝＝＝＝＝＝＝＝＝＝＝＝＝
\input{include_cv_pi}

% Co-investigators　＝＝＝＝＝＝＝＝＝＝＝＝＝＝＝
\input{include_cv_coi}

% To add co-investigators, make a copy of include_cv_coi.tex, rename it, and 
% add the file below.

%\input{include_cv_taro}
%\input{include_cv_hanako}


%#Split: 99_tail
\input{pieces/hook9} % pieces
\end{document}

